\frontpage{Abstract}
Spot spraying applies herbicide only where weeds grow instead of over the whole field, but it needs a
vision system that can tell a weed from a crop plant. Supervised networks do this well after training
on labelled images of the target field; zero-shot foundation models promise to do it with no training
at all. \textbf{AgriAgent} tests that promise on real field photographs. It implements the Grounded SAM
pipeline, a text-prompted Grounding DINO detector
followed by the SAM~2 segmenter, and adds four things around it: a deterministic
safety gate that is the only authority on what may be sprayed, a rule-based plant classifier built on
the excess-green index and the planting row, one bounded escalation agent
that consults vision-language models about uncertain plants, and herbicide advice taken only from the
ICAR herbicide bulletin. The whole pipeline runs on a laptop GPU (RTX~5050,
8~GB) in about 0.6~s per photo.

We evaluate on sugar beet and cotton photographs whose labels were copied from the official dataset
files, after discovering and replacing an earlier evaluation set with fabricated labels. On 90
held-out sugar beet photos, the Grounded SAM label policy sprays \textbf{65.5\% of crop pixels} while
reaching only \textbf{1.8\% of weed pixels}: text prompts alone do not separate crop from weed. Four
frozen rules (crop-mask overlap, planting row, size relative to the crop, off-row position) raise weed
recall to \textbf{77.9\%} and cut crop pixels sprayed to \textbf{8.8\%}, measured once on the held-out set,
while saving 95\% of the herbicide of blanket spraying. The escalation agent never approved spraying a
crop plant and was right in 13 of its 14 ``crop'' verdicts, but it called no weed a weed: it is safe
but not effective. These findings agree with the best published zero-shot crop--weed
work and show how much a transparent, testable safety layer recovers without
any training. AgriAgent is a research prototype and must not drive a real sprayer.

\vspace{1em}
\noindent\textbf{Keywords:} precision agriculture, weed detection, zero-shot segmentation, Grounding DINO,
SAM~2, safety gate, vision-language models, spot spraying.
\clearpage
